% =============================================================================
%  CHAPTER 7 -- CONCLUSION                                            [v4 OWNS]
% -----------------------------------------------------------------------------
%  v4.1 (2026-09-24): the Conclusion is Chapter 7 and holds the summary and the answers to
%  the research questions only (author); the deployment view and the future work are in
%  Chapter 8, the Discussion. Text: the v4.0 Summary and RQ answers, unchanged
%  (withheld/20260924_v4.1_archive/08_conclusion_v4.0.tex). Written on Chapters 5-6 at
%  v3.98 and Chapters 1-4 at v2.27 (via v3); every number restates one of Chapter 6.
%  Labels ch:conclusion, sec:conc:summary, sec:conc:rqs are kept.
%  v4.2 (2026-09-25): the ChatGPT audit of v4.1a applied (AUDIT_v4.1a_2026-09-25.md, section 15,
%  items C1-C8) on Chapters 5-6 at v3.100: CI-MeanFM dominates on D3IL-avoiding and MeanFM trades
%  one episode; the aligning claims scoped (unprojected K=20, medians, 1.8 cm test); the model
%  order scoped to the tasks with human demonstrations, the corridor tie to traversing cells, the
%  s-curve's own order named; \nfe for budgets; RQ2 without the 2.45x claim.
%  Individual changelog: changelogs/v4.2_20260925_audit_v4.1a_applied.md
% =============================================================================

\chapter{Conclusion}\label{ch:conclusion}

\section{Summary}\label{sec:conc:summary}

This thesis replaced the diffusion model of \ac{DPCC} with three flow-based objectives,
instantaneous-velocity matching, analytic average-velocity matching and consistency-interpolated
average-velocity matching, and set endpoint projection beside the per-step projection of \ac{DPCC},
keeping the temporal U-Net, the projection program, the data and the evaluation of the baseline. The
step budget thereby moved from training to inference. The comparison was made in three steps: on the
obstacle-avoidance benchmark of \ac{DPCC} with its setup unchanged, on a vision-conditioned alignment
task built for this thesis, and on a quadrotor benchmark of three scenes built for it as well.

On the benchmark of \ac{DPCC} consistency-interpolated average-velocity matching Pareto-dominates its
diffusion model with its own U-Net: the same success with constraint satisfaction, fewer control
steps, and about a thirtieth of the time per action, twelve times less than a baseline retrained at
two denoising steps; analytic average-velocity matching matches the saving at one episode in thirty.
Instantaneous-velocity matching dominates the baseline as well, so the saving belongs first to a
flow-based sampler at a small budget and then to the average-velocity objective, which adds the fewest
control steps. On D3IL-aligning the dominance repeats where the task separates the models: without projection, at
twenty sampling steps, analytic average-velocity matching is the only model that moves the box in
every context, and it reduces the median final distance by $84\,\%$ of the mean initial distance,
$0.0741$ against $0.4530$\,m, where the consistency-interpolated model, instantaneous-velocity
matching and the baseline leave the box near where it started; its median falls at every tested
budget, while the consistency-interpolated model's varies from two to twenty steps and falls
substantially only at a hundred, still behind. No context ends within the $1.8$\,cm of the task's own
position test, under the selected endpoint configuration at twenty steps or without projection.
Endpoint projection has no step to guide at the one-evaluation operating point of D3IL-avoiding and
nothing to win there; on D3IL-aligning, run at $\nfe = 20$, it is the better projector on the
constraint, keeping every context violation-free at the price of per-step projection.

The quadrotor scenes carry the framework to a plant that holds nothing by itself. UAV-pillars flies
the operating points of D3IL-avoiding and keeps their result on the declared constraints, and shows
what the plant changes: the vehicle tracks more closely than the arm, and a collision with a pillar the
constraints leave out ends the flight. UAV-corridor shows that the projection scheme of \ac{DPCC}
carries over to three dimensions, against constraints that leave the horizontal plane and that no
demonstration satisfies, and that on straight-line demonstrations the generative model does not
matter, the budget does; there the two projectors trade against each other without one dominating.
UAV-s-curve shows that the controller matters: the same planner configuration succeeds or fails by
the controller that flies it in closed loop, and the cascaded geometric controller used throughout is
the better of the two, at a twenty-fourth of the controller's cost per control step --- controller
and simulator step together, on the unprojected pair, a tenth of the whole loop.

With human demonstrations the average-velocity models lead, on the harder tasks the model and the
projection matter together and endpoint projection helps where the task needs a large budget; with
simple generated demonstrations instantaneous-velocity matching is enough: on the corridor the
objectives cannot be told apart where the flights traverse it at a jointly tested budget, and on the
s-curve instantaneous-velocity matching leads; and on the quadrotor the plant and its controller
matter as much as the planner.

\section{Answers to the Research Questions}\label{sec:conc:rqs}

\begin{description}
  \item[RQ1 --- Generative model.] Yes, with a margin, on the benchmark of \ac{DPCC}. With the
    backbone, the data and the evaluation held fixed, consistency-interpolated average-velocity
    matching and instantaneous-velocity matching reach the diffusion baseline's success with
    constraint satisfaction on D3IL-avoiding at one network evaluation, where the diffusion model,
    run away from its trained budget, loses it: $1.000$ in $59.2$ control steps at $18.1$\,ms per
    action and $1.000$ in $67.0$ at $17.3$, against the baseline's $1.000$ in $70.1$ at $553.4$;
    analytic average-velocity matching reads $0.967$ in $58.6$ at $18.3$. On D3IL-avoiding and
    D3IL-aligning, the tasks with human demonstrations, the order is the average-velocity models
    first, then instantaneous-velocity matching and the diffusion model, which are level on outcome
    and differ in cost. On D3IL-aligning without projection analytic average-velocity matching
    reduces the median final distance by $84\,\%$ of the mean initial distance at $\nfe = 20$,
    against $14$, $10$ and $-2\,\%$ for the consistency-interpolated model, instantaneous-velocity
    matching and the baseline, and under per-step projection it keeps nine contexts violation-free
    against the baseline's four at a third of its time. Between the two average-velocity objectives
    the consistency-interpolated one is ahead by one episode in thirty on D3IL-avoiding, keeps that
    lead on UAV-pillars, which flies both, and is behind by most of the task on D3IL-aligning, where
    its median varies from two to twenty steps and falls substantially only at a hundred, still
    behind the analytic objective, whose median falls at every tested budget. On the quadrotor's
    generated demonstrations the picture is the corridor's and the s-curve's: on the corridor the
    objectives cannot be told apart where the flights traverse it at a jointly tested budget, and on
    the s-curve's single route instantaneous-velocity matching crosses the finish line on nine
    flights of ten, the consistency-interpolated model on six and the analytic one on three, with the
    fewest violating steps. One evaluation saturates D3IL-avoiding, D3IL-aligning is run at its
    selected operating budget of twenty steps, and on the s-curve more sampling steps lose flights.
  \item[RQ2 --- Constraints.] Endpoint projection guides the sample only where the budget and the
    threshold leave a sampling step before the terminal one, the guiding-step count of
    \autoref{sec:method:degenerate}: at $\nfe = 1$ never, and at $\nfe = 2$ with the threshold of
    $0.5$ not, where it reduces to a terminal projection of the finished sample and gives no evidence
    of guidance; on instantaneous-velocity matching a run at $\nfe = 2$ that handed both methods the
    same candidate plans produced bit-identical rollouts, while on the two average-velocity models the
    two arms differ even there. Where it has steps to guide it is the better projector on the
    constraint at a comparable planning time: on D3IL-aligning at $\nfe = 20$ it keeps more contexts
    violation-free than per-step projection in seven of nine model--rule comparisons, on analytic
    average-velocity matching under random selection all ten against nine at $275$ against
    $266$\,ms per control step, and at $\nfe = 10$ under the same rule at $37\,\%$ of per-step
    projection's time; on D3IL-avoiding at $\nfe = 3$ it reaches $1.000$ on the two average-velocity
    models where per-step projection reaches $0.917$ and $0.958$, at about half the time per control
    step. It does not reach the goal in fewer control steps, and on UAV-corridor it trades: fewer and
    shallower violations at $\nfe = 20$, cheaper over the hump with one candidate and dearer under
    the tilt, on longer flights. It helps on the hard task that needs the budget, D3IL-aligning, and
    not on D3IL-avoiding, which is solved at a budget where it has no step to act.
  \item[RQ3 --- Transfer.] The answers hold where the demonstrations are of the same kind, and the
    plant and the demonstrations change the rest. Analytic average-velocity matching leads again on
    camera observations and contact, and the declared constraints transfer exactly from the
    manipulator to the quadrotor on UAV-pillars, where the two average-velocity models keep their
    order. The quadrotor adds what the manipulator could not show: an undeclared obstacle ends a
    flight, the controller decides whether a plan on a route of sharp turns is flown at all, and on
    the corridor's demonstrations, one path forward from every position, the objectives cannot be
    told apart where the flights traverse it. The projection scheme itself carries over to three
    dimensions and to constraints that leave
    the horizontal plane.
\end{description}
\dataref{Every number restates \autoref{sec:res:summary}: \autoref{tab:avoiding-conclusion},
\autoref{tab:va-models}, \autoref{tab:va-projection-models}, \autoref{tab:avoiding-projectors},
\autoref{sec:res:uav:projection}, \autoref{tab:uav-pillars-raw}, \autoref{tab:uav-controller}.}
